%% kapa: exam=2 round=22 year=2011 period=1 subject=appraisal_practice subject_ko=감정평가실무 seq=1 number=1 points=40
%% source: Q-Net 감정평가사 기출문제 (cst00309) artlSeq=5040389 file=extracted/1008/5040389/2011년 제22회 감정평가사 2차시험 기출문제/1교시.hwp
%% conversion: kordoc 4.24.51 -> tools/gen2.py (auto); review=proofread-sonnet; flags=none
\begin{longquestion}{【문제 1】}
감정평가사 甲은 (주)K생명보험으로부터 동 회사가 보유 중인 부동산에 대한 감정평가를 의뢰받아 처리계획을 수립한 후 현장조사를 수행하고 아래와 같이 자료를 수집ㆍ분석하였는바, 이를 활용하여 시산가치 조정을 통한 최종 감정평가액을 산출하되, 평가방식 적용 시 필요한 경우 그 판단에 대한 의견을 명기하시오. (40점) \par
<자료 1> 기본적 사항 \par
- 감정평가 의뢰 내역 \par
{\small\centering
\begin{tabular}{|c|c|c|c|c|c|}
\hline
\makecell{항목} & \makecell{기호} & \makecell{소재지} & \makecell{지목\newline{}층} & \makecell{면적(m\textsuperscript{2})} & \makecell{용도지역\newline{}용도} \\
\hline
\multirow{6}{*}{\makecell{토지}} & \makecell{1} & \makecell{경기도 Y시 B읍 K리 219} & \makecell{잡종지} & \makecell{225.0} & \makecell{계획관리} \\
\cline{2-6}
 & \makecell{2} & \makecell{〃 219-2} & \makecell{잡종지} & \makecell{291.0} & \makecell{〃} \\
\cline{2-6}
 & \makecell{3} & \makecell{〃 219-1} & \makecell{대} & \makecell{975.0} & \makecell{〃} \\
\cline{2-6}
 & \makecell{4} & \makecell{〃 219-3} & \makecell{대} & \makecell{554.0} & \makecell{〃} \\
\cline{2-6}
 & \makecell{5} & \makecell{〃 219-4} & \makecell{도로} & \makecell{105.0} & \makecell{〃} \\
\cline{2-6}
 & \makecell{6} & \makecell{〃 219-5} & \makecell{도로} & \makecell{144.0} & \makecell{〃} \\
\hline
\multirow{2}{*}{\makecell{건물}} & \makecell{가} & \makecell{〃 219-1\newline{}〃 219-3} & \makecell{지하1\newline{}지상3} & \makecell{1,254.3} & \makecell{근린생활시설\newline{}및 숙박시설} \\
\cline{2-6}
 & \makecell{나} & \makecell{〃 219-1\newline{}〃 219-3} & \makecell{1} & \makecell{72.24} & \makecell{근린생활시설} \\
\hline
\end{tabular}\par}

- 의 뢰 인 : (주)K생명보험 \par
- 평가목적 : 일반거래(매매참고용) \par
- 제 출 처 : (주)K생명보험 \par
- 목록표시근거 : 등기부등본, 일반건축물관리대장등본 \par
- 가격시점 : 2011. 09. 04 \par
- 조사기간 : 2011. 08. 31 ～ 2011. 09. 04 \par
- 작성일자 : 2011. 09. 04 \par
<자료 2> 대상부동산에 관한 기본 자료 및 현장조사 내용 \par
1) 토지 \par
- 이용상황 \par
ㆍ기호 (1, 2) : 본건 기호 (5, 6) 토지와 경계 구분 없이 본건 (가, 나) 건물의 진출입을 위한 포장도로로 이용 중임 \par
ㆍ기호 (3) : 현황 숙박시설로 이용 중인 본건 기호 (가) 건물 부지임 \par
ㆍ기호 (4) : 현황 주택으로 이용 중인 본건 기호 (나) 건물 부지임 \par
ㆍ기호 (5, 6) : 본건 기호 (1, 2) 토지와 경계 구분 없이 본건 (가, 나) 건물 진출입을 위한 포장도로로 이용 중임 \par
- 접면도로 \par
ㆍ본건 기호 (2, 6) 토지가 남서측으로 ○○번 국도에 접함 \par
- 형상 및 지세 \par
ㆍ본건 일단의 토지 전체를 기준으로 자루형에 가까운 평지임 \par
2) 건물 \par
- 공통사항 \par
ㆍ사용승인일 : 1997. 08. 29 \par
ㆍ건폐율 / 용적률 : 24.91\% / 65.07\% \par
- 기호 (가) \par
ㆍ구조 : 철근콘크리트 슬라브지붕 \par
ㆍ난방설비, 패키지 에어컨에 의한 냉방설비, 위생 및 급배수설비, 화재탐지ㆍ경보 및 소화설비 등 \par
ㆍ이용상황 \par
{\small\centering
\begin{tabular}{|c|c|c|c|}
\hline
\makecell{층} & \makecell{용도} & \makecell{면적} & \makecell{이용상황} \\
\hline
\makecell{지1} & \makecell{근린생활시설} & \makecell{331.55} & \makecell{창고, 보일러실} \\
\hline
\makecell{1} & \makecell{숙박시설} & \makecell{305.23} & \makecell{접수대, 객실 9} \\
\hline
\makecell{2} & \makecell{숙박시설} & \makecell{308.76} & \makecell{객실 10} \\
\hline
\makecell{3} & \makecell{숙박시설} & \makecell{308.76} & \makecell{객실 10} \\
\hline
\end{tabular}\par}

- 기호 (나) \par
ㆍ구조 : 연와조 슬라브지붕 단층 \par
ㆍ면적 : 72.24m\textsuperscript{2} \par
ㆍ주요설비 : 난방설비, 위생 및 급배수설비 등 \par
ㆍ이용상황 : 주택(방2, 화장실 겸 욕실, 주방 겸 거실) \par
3) 임대 및 사용현황 \par
- 2009년부터 乙에게 무상으로 임대 중인 것으로 조사되었음 \par
4) 본건 평가전례 \par
- 가격시점 : 2006. 03. 02 \par
- 평가목적 : 담보 \par
- 평 가 액 : 1,402,384,500원 \par
<자료 3> 경기도 Y시 개황 \par
1) Y시는 경기도 북동측 내륙권, 서울에서 직선거리 약 30km 거리에 위치하며 대체로 산지가 많고 평지가 적음 \par
2) 도로 및 전철 등 서울로의 대중교통 접근조건이 개선됨에 따라 대규모 택지개발사업 시행 및 그로 인하여 인구유입이 지속적으로 증가되어 왔고, 이에 따라 주택이 지속적으로 공급되어 왔음 \par
3) Y시의 산업별 사업체 수는 제조업, 도소매업, 숙박(음식점)업 순으로 나타났고, 종사인원을 기준으로 할 경우 제조업이 50\%정도를 차지하나 대형 할인점 등의 입지에 따른 3차 산업의 발달 등 급격한 도시화로 도소매업의 증가가 두드러진 반면 숙박업은 지속적인 감소세를 나타내고 있음 \par
4) Y시 토지거래는 2000년대 초반, 시 승격을 전후로 연간 6,500여건에서 연간 16,000여건으로 거래량이 급격히 상승하였으나 2000년대 중반 이후 등락을 보이다 최근에는 연간 10,000여건으로 안정세를 보이고 있음 \par
<자료 4> 지역 및 대상부동산 개황 \par
1) 본건은 Y시 B읍 K리 소재 K저수지 북동측 인근에 위치하고 있으며, 인근은 국도 및 지방도변을 따라 숙박시설 및 근린생활시설 등이 산재하고 후면으로 농경지 및 임야 등이 혼재하여 있음 \par
2) 본건 토지 남서측으로 ○○번 국도에 접하고 있어 본건까지 제반 차량의 진출입이 원활하고, 본건 북서측 인근에 ○○번 국도와 동서로 연결되는 □□번 지방도가 연결되는 삼거리가 소재함 \par
3) Y시 관내를 연결하는 버스정류장이 본건 남동측 인근에 위치하나, 운행 간격 등에 비추어 대중교통 사정은 다소 불편함 \par
<자료 5> 인근지역 분석 \par
1) 본건은 계획관리지역 내 3층 숙박시설(객실 29개)로서, 용도적ㆍ기능적 동일성을 기준으로 본건 남서측 인근에 소재한 K저수지 북측의 ○○번 국도변 및 동 국도에 동서로 연결되는 □□번 지방도를 따라 본건 북서측 인근 M저수지에 이르는 지방도변 일대가 본건 부동산의 인근지역으로 판단됨 \par
2) 인근지역은 K저수지를 중심으로 유원지ㆍ낚시터 등의 이용객들을 위한 음식점, 숙박시설, 카페 등이 산재되어 있고, M저수지 주변도 이와 유사한 이용을 보이고 있으며 특히, 숙박시설은 독립적으로 위치하기 보다는 몇 개씩 집단화고 있는 양상을 보이고 있으나, 경기침체 및 유원지ㆍ낚시터 등의 이용객의 감소와 더불어 숙박시설이 집단화한 지역으로서의 전반적인 경쟁력 약화 등이 상승 작용하여 영업상황이 악화되어 가고 있으며, 영업을 중지하는 숙박시설이 증가하고 있는 추세임 \par % [원문 그대로] 집단화고 있는
3) 인근지역은 ○○번 국도 및 □□번 지방도를 따라 노변 또는 후면에 음식점 또는 숙박시설이 주류를 이루고 있으며, 숙박시설의 경우 3층으로서 객실 30～35개 규모가 일반적 이용임 \par
4) 토지 가격수준에 대한 탐문 조사결과, ○○번 국도에 접하고 있는 경우 @453,000원/m\textsuperscript{2}～@600,000원/m\textsuperscript{2}, 후면지의 경우 @272,000원/m\textsuperscript{2}～@300,000원/m\textsuperscript{2} 수준이고, M저수지에서 동측으로 □□번 지방도에 접하고 있는 경우 @453,000원/m\textsuperscript{2}～@544,000원/m\textsuperscript{2}, 후면지의 경우 @211,000원/m\textsuperscript{2}～@300,000원/m\textsuperscript{2} 수준으로 호가되고 있음 \par
5) 또한, 최근 인근지역에는 영업이 중단된 숙박시설을 노인전문요양원 등 타 용도로 전환ㆍ이용하려는 목적으로 매물을 찾는 문의가 간혹 있으나 실제 노인전문요양원 등으로 전환ㆍ이용된 사례는 없는 것으로 조사되었음 \par
<자료 6> 공시지가 표준지 및 매매사례 \par
1) 비교 표준지 공시지가 \par
{\small\centering
\setlength{\tabcolsep}{3pt}
\begin{tabular}{|c|c|c|c|c|c|c|c|c|c|}
\hline
\makecell{소재지} & \makecell{면적\newline{}(m\textsuperscript{2})} & \makecell{지목} & \makecell{이용\newline{}상황} & \makecell{용도\newline{}지역} & \makecell{도로\newline{}교통} & \makecell{형상\newline{}지세} & \makecell{공시기준일} & \makecell{공시지가} & \makecell{비고} \\
\hline
\multirow{5}{*}{\makecell{K리\newline{}217-1}} & \multirow{5}{*}{\makecell{880.0}} & \multirow{5}{*}{\makecell{대}} & \multirow{5}{*}{\makecell{상업용}} & \multirow{5}{*}{\makecell{계획\newline{}관리}} & \multirow{5}{*}{\makecell{소로\newline{}각지}} & \multirow{5}{*}{\makecell{부정형\newline{}평지}} & \makecell{2007.01.01} & \makecell{405,000} & \multirow{5}{*}{\parbox{2cm}{\centering 본건 남동측 인접, ○○번 국도변}} \\
\cline{8-9}
 &  &  &  &  &  &  & \makecell{2008.01.01} & \makecell{455,000} &  \\
\cline{8-9}
 &  &  &  &  &  &  & \makecell{2009.01.01} & \makecell{420,000} &  \\
\cline{8-9}
 &  &  &  &  &  &  & \makecell{2010.01.01} & \makecell{425,000} &  \\
\cline{8-9}
 &  &  &  &  &  &  & \makecell{2011.01.01} & \makecell{430,000} &  \\
\hline
\end{tabular}\par}

2) 매매사례 \par
{\small\centering
\setlength{\tabcolsep}{2pt}
\begin{tabular}{|c|c|c|c|c|c|}
\hline
\multicolumn{2}{|c|}{\makecell{구분}} & \makecell{매매사례 \#1} & \makecell{매매사례 \#2} & \makecell{매매사례 \#3} & \makecell{매매사례 \#4} \\
\hline
\multicolumn{2}{|c|}{\makecell{소재지}} & \makecell{K리 354} & \makecell{K리 419-3} & \makecell{K리 418} & \makecell{K리 381-5} \\
\hline
\multicolumn{2}{|c|}{\makecell{매매금액}} & \makecell{481,100,000} & \makecell{903,500,000} & \makecell{685,100,000} & \makecell{850,100,000} \\
\hline
\multicolumn{2}{|c|}{\makecell{매매일자}} & \makecell{2009.06.05} & \makecell{2010.01.07} & \makecell{2009.08.20} & \makecell{2011.02.27} \\
\hline
\multicolumn{2}{|c|}{\makecell{토지면적(m\textsuperscript{2})}} & \makecell{974.0} & \makecell{1,327.0} & \makecell{1,405.0} & \makecell{1,258.0} \\
\hline
\multirow{6}{*}{\makecell{건물\newline{}내역}} & \makecell{구조} & \makecell{-} & \makecell{철근콘크리트조\newline{}슬라브지붕} & \makecell{철근콘크리트조\newline{}슬라브지붕} & \makecell{철근콘크리트조\newline{}슬라브지붕} \\
\cline{2-6}
 & \makecell{용도} & \makecell{-} & \makecell{숙박시설} & \makecell{숙박시설} & \makecell{숙박시설} \\
\cline{2-6}
 & \makecell{객실 수} & \makecell{-} & \makecell{31} & \makecell{26} & \makecell{34} \\
\cline{2-6}
 & \makecell{층\newline{}(지상/지하)} & \makecell{-} & \makecell{3/1} & \makecell{5/1} & \makecell{3/1} \\
\cline{2-6}
 & \makecell{연면적(m\textsuperscript{2})} & \makecell{-} & \makecell{1,349.74} & \makecell{975.24} & \makecell{2,410.27} \\
\cline{2-6}
 & \makecell{사용승인일} & \makecell{-} & \makecell{2004.12.24} & \makecell{1999.08.15} & \makecell{2002.09.01} \\
\hline
\multicolumn{2}{|c|}{\makecell{비고}} & \makecell{-} & \makecell{-} & \makecell{현 영업중단} & \makecell{현 영업중단} \\
\hline
\end{tabular}\par}

<자료 7> 표준건축비 자료(가격시점 기준) \par
1) 철근콘크리트조 숙박시설 : @1,060,000원/m\textsuperscript{2} \par
2) 연와조 단독주택 : @850,000원/m\textsuperscript{2} \par
<자료 8> 시점수정 자료 \par
- 2007. 01. 01. 이후 인근지역 토지시장 및 건축물 신축가격은 큰 변화가 없었던 것으로 조사되어 시점수정은 필요 없음 \par
<자료 9> 본건 영업자료 등 \par
1) 매출액 : 5,500,000원/월 \par
2) 영업경비 \par
- 단기내용연수 항목에 대한 대체충당금을 제외한 세금 등 고정비와 인건비, 냉난방비, 공과금 등 변동비에 대한 자료검토 결과 평균 4,500,000원/월 소요됨 \par
3) 동산항목 가치 \par
- 가격시점 현재 본건 숙박시설 내 가구, 전자제품 등의 잔존가치는 객실당 @600,000원으로 산정됨 \par
<자료 10> 요인비교 자료 \par
1) 토지 개별요인 비교치 \par
{\small\centering
\begin{tabular}{|c|c|c|c|c|c|}
\hline
\makecell{공시지가 표준지} & \makecell{본건\newline{}(전체 기준)} & \makecell{매매사례 \#1} & \makecell{매매사례 \#2} & \makecell{매매사례 \#3} & \makecell{매매사례 \#4} \\
\hline
\makecell{1.00} & \makecell{0.78} & \makecell{1.15} & \makecell{1.07} & \makecell{0.97} & \makecell{0.95} \\
\hline
\end{tabular}\par}

2) 건물 개별요인 비교치 \par
- 철근콘크리조 \par
{\small\centering
\begin{tabular}{|c|c|c|c|c|}
\hline
\makecell{표준건축비} & \makecell{본건} & \makecell{매매사례 \#2} & \makecell{매매사례 \#3} & \makecell{매매사례 \#4} \\
\hline
\makecell{1.00} & \makecell{1.00} & \makecell{1.00} & \makecell{1.03} & \makecell{1.00} \\
\hline
\end{tabular}\par}

- 연와조 \par
{\small\centering
\begin{tabular}{|c|c|}
\hline
\makecell{표준건축비} & \makecell{본건} \\
\hline
\makecell{1.00} & \makecell{1.05} \\
\hline
\end{tabular}\par}

<자료 11> 각종 이율 등 \par
1) 시장금리 등 \par
- 국고채(3년) : 3.96\%/연 \par
- 회사채(3년,AA-) : 6.12\%/연 \par
- 저당이자율(예금은행 가중평균 대출금리) : 6.73\%/연 \par
2) 숙박시설 지분환원율 등 \par
- 지분환원율 : 12.0\%/연 \par
- 대출비율 : 45\% \par
3) 숙박시설 투자수익률 등 \par
- 투자수익률 : 15.0\%/연 \par
- 가치변동 : 전형적인 보유기간인 5년 동안 인근지역 숙박시설의 가치는 35\% 정도 하락할 것으로 예상 \par
<자료 12> 기타 \par
1) 원가방식 중 내용연수법 적용 시 내용연수 또는 잔존내용연수 조정이 필요할 경우 매매사례를 활용 \par
2) 비교방식 중 개별요인 비교는 사례의 토지/건물 가치구성비율을 활용하고, 내용연수 만료 시 잔존가치율은 1\%를 적용 \par
3) 숙박시설 운영을 위해서는 직원 등의 숙소가 필요함이 일반적인 바, 본건 기호 (나) 건물은 무상임차인 乙과 운영을 돕고 있는 그 자녀 1인의 숙소로 사용되는 바, 숙박시설 용도에 부합하는 것으로 판단 \par
4) 가격시점 현재 건축물의 철거비는 @15,000원/m\textsuperscript{2} \par
\end{longquestion}
